Im $V(T)$-Diagramm ist eine Isobare eine Gerade durch den Ursprung ($V \propto T$), eine Isotherme senkrecht und eine Isochore waagrecht. Der Druck folgt aus $p/p_0 = \frac{T/T_0}{V/V_0}$.
\begin{abcliste}
\abc A $\to$ B ist $\result{\text{isobar}}$ (Expansion, $V$ und $T$ verdreifachen sich), B $\to$ C $\result{\text{isotherm}}$ (Kompression bei $T = 3\,T_0$), C $\to$ A $\result{\text{isochor}}$ (Abkühlung bei $V = V_0$).
\begin{center}
\begin{tabular}{lccc}
 & A & B & C \\ \hline
 $T$ in $T_0$ & 1 & 3 & 3 \\
 $V$ in $V_0$ & 1 & 3 & 1 \\
 $p$ in $p_0$ & 1 & 1 & 3
\end{tabular}
\end{center}
Also $\result{p_A = p_B = p_0}$, $\result{p_C = 3\,p_0}$.

\abc Im $p(V)$-Diagramm ist die Isobare A $\to$ B waagrecht, die Isotherme B $\to$ C eine Hyperbel ($pV = 3\,p_0V_0$) und die Isochore C $\to$ A senkrecht:
\begin{center}
\begin{tikzpicture}[x=1cm, y=1cm, >=Stealth, font=\small]
 \draw[black!12] (0,0) grid[xstep=1, ystep=1] (4.001,4.001);
 \draw[->] (0,0) -- (4.6,0) node[right] {$V$ in $V_0$};
 \draw[->] (0,0) -- (0,4.6) node[above] {$p$ in $p_0$};
 \node[below left] at (0,0) {\small 0};
 \draw (1,0.08) -- (1,-0.08) node[below] {\small 1};
 \draw (2,0.08) -- (2,-0.08) node[below] {\small 2};
 \draw (3,0.08) -- (3,-0.08) node[below] {\small 3};
 \draw (4,0.08) -- (4,-0.08) node[below] {\small 4};
 \draw (0.08,1) -- (-0.08,1) node[left] {\small 1};
 \draw (0.08,2) -- (-0.08,2) node[left] {\small 2};
 \draw (0.08,3) -- (-0.08,3) node[left] {\small 3};
 \draw (0.08,4) -- (-0.08,4) node[left] {\small 4};
 \draw[very thick, blue] (1,1) -- (3,1);
 \draw[-{Stealth[length=2.6mm]}, blue] (2.09,1) -- (2.11,1);
 \draw[very thick, blue, domain=1:3, samples=60, smooth] plot (\x, {3/\x});
 \draw[-{Stealth[length=2.6mm]}, blue] (1.739,1.725) -- (1.725,1.739);
 \draw[very thick, blue] (1,3) -- (1,1);
 \draw[-{Stealth[length=2.6mm]}, blue] (1,1.91) -- (1,1.89);
 \fill[black] (1,1) circle (2pt) node[below left] {$A$};
 \fill[black] (3,1) circle (2pt) node[below right] {$B$};
 \fill[black] (1,3) circle (2pt) node[above left] {$C$};
 \draw[black!35, dashed, domain=0.8:4, samples=40, smooth] plot (\x, {3/\x});
\end{tikzpicture}
\end{center}

\abc Im $p(T)$-Diagramm ist die Isobare A $\to$ B waagrecht, die Isotherme B $\to$ C senkrecht und die Isochore C $\to$ A eine Gerade durch den Ursprung:
\begin{center}
\begin{tikzpicture}[x=1cm, y=1cm, >=Stealth, font=\small]
 \draw[black!12] (0,0) grid[xstep=1, ystep=1] (4.001,4.001);
 \draw[->] (0,0) -- (4.6,0) node[right] {$T$ in $T_0$};
 \draw[->] (0,0) -- (0,4.6) node[above] {$p$ in $p_0$};
 \node[below left] at (0,0) {\small 0};
 \draw (1,0.08) -- (1,-0.08) node[below] {\small 1};
 \draw (2,0.08) -- (2,-0.08) node[below] {\small 2};
 \draw (3,0.08) -- (3,-0.08) node[below] {\small 3};
 \draw (4,0.08) -- (4,-0.08) node[below] {\small 4};
 \draw (0.08,1) -- (-0.08,1) node[left] {\small 1};
 \draw (0.08,2) -- (-0.08,2) node[left] {\small 2};
 \draw (0.08,3) -- (-0.08,3) node[left] {\small 3};
 \draw (0.08,4) -- (-0.08,4) node[left] {\small 4};
 \draw[black!30, dashed] (0,0) -- (1,1);
 \draw[very thick, blue] (1,1) -- (3,1);
 \draw[-{Stealth[length=2.6mm]}, blue] (2.09,1) -- (2.11,1);
 \draw[very thick, blue] (3,1) -- (3,3);
 \draw[-{Stealth[length=2.6mm]}, blue] (3,2.09) -- (3,2.11);
 \draw[very thick, blue] (3,3) -- (1,1);
 \draw[-{Stealth[length=2.6mm]}, blue] (1.907,1.907) -- (1.893,1.893);
 \fill[black] (1,1) circle (2pt) node[above left] {$A$};
 \fill[black] (3,1) circle (2pt) node[below right] {$B$};
 \fill[black] (3,3) circle (2pt) node[above right] {$C$};
\end{tikzpicture}
\end{center}

\abc Jeder Schritt behält seine Art, aber nicht seine Form: Die Isotherme B $\to$ C ist im $V(T)$-Diagramm senkrecht, im $p(V)$-Diagramm aber eine $\result{\text{Hyperbel}}$. Ein Dreieck aus drei Geraden kann es im $p(V)$-Diagramm nicht geben. Zudem müsste dort die Isochore senkrecht und die Isobare waagrecht sein, nicht schräg.
\end{abcliste}
